% Copyright (c) 2026 Twin Vector Labs LLC.
% All rights reserved.
\documentclass[11pt]{article}

\usepackage[margin=1in]{geometry}
\usepackage{amsmath,amssymb,amsthm}
\usepackage{mathtools}
\usepackage{enumitem}
\usepackage{booktabs}
\usepackage{array}
\usepackage[colorlinks=true,linkcolor=blue,urlcolor=blue]{hyperref}

\newtheorem{theorem}{Theorem}
\newtheorem{proposition}{Proposition}
\newtheorem{lemma}{Lemma}
\theoremstyle{definition}
\newtheorem{definition}{Definition}
\newtheorem{remark}{Remark}

\newcommand{\C}{\mathbb{C}}
\newcommand{\R}{\mathbb{R}}
\newcommand{\Z}{\mathbb{Z}}
\newcommand{\rep}[1]{\mathbf{#1}}
\newcommand{\abs}[1]{\left\lvert #1 \right\rvert}
\newcommand{\norm}[1]{\left\lVert #1 \right\rVert}
\newcommand{\sigmacol}{\sigma}
\newcommand{\Pin}{P_{\mathrm{in}}}
\newcommand{\Pout}{P_{\mathrm{out}}}
\newcommand{\winA}{W_{\!A}}
\newcommand{\winB}{W_{\!B}}
\newcommand{\winC}{W_{\!C}}
\newcommand{\winR}{W_{\!R}}

\title{\bfseries The Proton Color Singlet from Symmetric Register Windows:\\
how tetrahedral $A_4$ symmetry intertwines the color $\Z_3$ at the $W_{ABC}$ junction}
\author{Tessera --- cobordism subsystem (\#398, completing \#396)}
\date{}

\begin{document}
\maketitle

\begin{abstract}
The trivalent junction $W_{ABC}$ (\#396) carries the $\Z_3$ shadow of the
three-quark color singlet $\varepsilon_{ijk}$ that the bipartite merge cannot
(\#382): three input ``windows'' $A,B,C$ and one emergent result window $R$, each a
triple of holes, all on one connected $2$-frequency geodesic icosahedron ($S^2$,
$42$ vertices, $80$ triangles) extruded $\times I$. After \#396 the singlet was
\emph{reachable} --- the input$\to$result transport is full-rank with the singlet in
its image --- but a \emph{natural} quark input reached only $\sim 0.74$: the greedily
placed windows are geometrically \emph{inequivalent}, so the per-window transport
blocks are unrelated matrices. This note records the principled fix and a finding
that was not anticipated. We place the four windows as \emph{one orbit of a
tetrahedral subgroup} $A_4$ of the icosahedral rotation group, seated at a system of
imprimitivity on the icosahedron's vertices. We then show that symmetric windows are
\emph{necessary but not sufficient}: the singlet emerges only when three conditions
hold together --- (i) a symmetry-respecting (uniform) metric, (ii) the input in the
$\omega$-\emph{representation} of the window-cycling symmetry $g$, and (iii) a
$g$-coherent signed read-out. Given these, the transport $M$ \emph{intertwines} the
color $\Z_3$, $M\Pin=\Pout M$, and the $\omega$-representation input transports to the
singlet line with overlap $0.999$ --- the proton, forced by symmetry and never
imposed. The residual $\sim 4\times 10^{-2}$ in the intertwining is a fixed-resolution
triangulation artifact (the $\times I$ extrusion is not $g$-symmetric even though the
base surface and metric are); it is metric-independent, and $0.999$ rather than an
exact $1$ is its signature. All numbers are measured through the production C++
\texttt{TransportCobordism}.
\end{abstract}

\tableofcontents

\section{Where we were}
\label{sec:context}

The proton-from-quarks program builds the baryon color singlet by relaxing a
Lorentzian Regge cobordism with pinned boundary color states and reading the
\emph{emergent} result off the relaxed geometry --- nothing is imposed. We write the
\emph{color charge} of a window $W$ for the net signed period it carries,
\begin{equation}
  \sigmacol_W \;=\; \sum_{\text{holes } h\in W} (\text{signed induced period of } h),
  \label{eq:charge}
\end{equation}
the discrete analogue of $\sum_i \psi_i$ on a color triple. Two prior results frame
the present one.

\paragraph{The bipartite obstruction (\#382).} Composing pairwise merges never
reaches the singlet: the emergent result stays color-charged
($\abs{\sigmacol}\approx 0.67$ for the tripartite sequence) at every step. This is
representation theory, not a numerical defect: the color singlet $\varepsilon_{ijk}$
is the three-index alternating invariant, absent from
$\rep{3}\otimes\rep{3}=\rep{6}\oplus\rep{\bar 3}$ and first appearing in
$\rep{3}\otimes\rep{3}\otimes\rep{3}$. A baryon from fundamentals is irreducibly an
$n=3$ object.

\paragraph{The tripartite junction (\#396).} The simplicial resolution is the
trivalent junction $W_{ABC}$: one connected base surface --- a $2$-frequency geodesic
icosahedron ($S^2$, $42$ vertices, $80$ triangles) --- minus $12$ vertex-disjoint
hole triangles grouped into four \emph{windows} of three, $A,B,C$ (the quark inputs)
and $R$ (the emergent result), extruded $\times I$ into a $3$-complex
(\texttt{prismCells}). The $12$ hole cycles satisfy exactly one relation --- the
global Stokes constraint below --- so $b_1 = 12-1 = 11$; the inputs are
\emph{independent} cycles (they do not average, escaping the bipartite obstruction),
and that single relation is junction charge conservation:
\begin{equation}
  \sum_{\text{all holes}} (\text{induced period}) = 0
  \quad\Longrightarrow\quad
  \sigmacol_R = -\,(\sigmacol_A+\sigmacol_B+\sigmacol_C).
  \label{eq:stokes}
\end{equation}
Confinement is then a property of \eqref{eq:stokes}: the four window charges must sum
to zero, so an isolated colored input cannot close on the junction; only the
color-neutral (singlet) combination does. The bipartite $\abs{\sigmacol}\approx 0.67$
residual is precisely such a non-confined, charged failure.

\paragraph{The open problem.} After \#396 the read-out is a linear map
$M:\C^9\to\C^3$ (nine input color amplitudes, three result amplitudes); it is
full-rank ($\operatorname{rank}M=3$) and the singlet lies in its image, so the
singlet is \emph{reachable}. Yet the natural neutral input (three identical
color-neutral pairs, read with each hole's own orientation) gives overlap
$\sim 0.74$. The diagnosis: with the windows placed \emph{greedily} (first twelve
vertex-disjoint triangles in face order) the per-window transport blocks
$M_A,M_B,M_C$ (the three $3\times 3$ sub-blocks of $M$) are geometrically
inequivalent --- $M_A^{-1}M_B$ is a general matrix, not a signed permutation. The
input that hits the singlet is therefore geometry-specific and unnatural. The fix is
to make the windows \emph{equivalent}, which is the subject of this note.

\section{Symmetric placement: the windows as an \texorpdfstring{$A_4$}{A4} orbit}
\label{sec:placement}

\subsection{The symmetry available}
The regular icosahedron has rotation group $I\cong A_5$ of order $60$; the
$2$-frequency geodesic subdivision inherits it as a group of simplicial
automorphisms permuting the $80$ sub-triangles and fixing the metric of the round
sphere. $A_5$ contains tetrahedral subgroups $T\cong A_4$ of order $12$. Fix one such
$A_4$. Because $\gcd(\abs{A_4},5)=1$ and the vertex stabilizer in $A_5$ is the
$5$-fold rotation $C_5$, $A_4$ acts \emph{freely and transitively} on the $12$
vertices (one orbit of $12$). It nevertheless admits a \emph{system of imprimitivity}
with four blocks of three,
\begin{equation}
  \{2,8,10\},\quad \{1,4,7\},\quad \{0,6,9\},\quad \{3,5,11\},
  \label{eq:blocks}
\end{equation}
which partition the twelve vertices; $A_4$ permutes the four blocks by its standard
transitive degree-$4$ action, and each block is fixed setwise by one of $A_4$'s four
$C_3$ subgroups (its three-fold axis).

\subsection{The construction}
A \emph{window} is the triple of vertex-disjoint corner sub-triangles seated at one
block of \eqref{eq:blocks} and cycled by that block's three-fold axis. Concretely the
corner sub-triangle at icosahedron vertex $v$ is $(v,\,m(v,p),\,m(v,q))$, where
$m(\cdot,\cdot)$ is the edge-midpoint vertex and $p,q$ are two neighbours of $v$; the
$C_3$ rotation about the block's axis cycles the three corners, hence the three holes.
The four windows $\winA,\winB,\winC,\winR$ are therefore a \emph{single} $A_4$ orbit
(their stabilizer is the block's own $C_3$, and $\abs{A_4}/3 = 4$). All $12$ holes are
vertex-disjoint (the $36$ vertices are distinct), and the topology is unchanged:
$b_1=11$ and the dual complex passes the manifold gate.

\begin{proposition}[Symmetric windows]
\label{prop:windows}
The four windows in the cobordism's own vertex labels are
\[
\begin{array}{llll}
\winA: & (2,13,14)\ (8,31,32)\ (10,22,36) &\quad \winB: & (1,23,24)\ (4,19,34)\ (7,30,39)\\[2pt]
\winC: & (0,16,18)\ (6,26,27)\ (9,33,41) &\quad \winR: & (3,15,29)\ (5,20,21)\ (11,37,38)
\end{array}
\]
each window a genuine $C_3$ orbit, all real faces of the base surface, the $36$
vertices distinct.
\end{proposition}

\paragraph{Generated, not hardcoded.} The windows are produced \emph{from the
symmetry} at build time: four $C_3$ generators (as $12$-vertex permutations), one
seed corner per window, and the midpoint lift (a base vertex maps by the permutation;
a midpoint $m(p,q)$ maps to $m(p',q')$). This is correct in whatever midpoint
numbering the build produces. The point is not pedantic: two independent attempts to
\emph{hardcode} the twelve triples both produced tables that were wrong against the
actual numbering. Generation is correct-by-construction and is guarded by build-time
assertions (each window a $C_3$ orbit, all holes real faces, $36$ distinct vertices).
The explicit data are in Appendix~\ref{app:data}.

\section{Necessary but not sufficient: three ingredients}
\label{sec:ingredients}

The central empirical finding is that symmetric \emph{windows alone do not produce the
singlet}. Three conditions must hold together --- a symmetry-respecting metric, the
input in the right representation, and a coherent read-out --- and all three are
phrased in terms of the same window-cycling symmetry, which we name first.

\begin{definition}[The window-cycling symmetry $g$]
\label{def:g}
$A_4$ acting on the four windows is the alternating group on four points; its
three-cycles fix one window and cycle the other three. Let $g\in A_4$ be the element
with
\begin{equation}
  g:\ \winR \mapsto \winR,\qquad \winA\mapsto\winB\mapsto\winC\mapsto\winA .
  \label{eq:g}
\end{equation}
Such $g$ exists and is unique up to its inverse; on the result window it acts as a
$C_3$ on $R$'s three holes. \emph{This $g$ is the color $\Z_3$.}
\end{definition}

We write $\Pin$ (resp.\ $\Pout$) for the signed permutation by which $g$ acts on the
nine input-hole (resp.\ three result-hole) period amplitudes; both are made precise in
\S\ref{sec:theorem}.

\subsection{(i) A symmetry-respecting metric}
The seed metric must be invariant under $g$. A random ``jitter'' seed --- previously
used to lift the degenerate uniform point --- breaks the $A_4$ symmetry: charge
conservation \eqref{eq:stokes} degrades from $\sim 10^{-14}$ to $\sim 0.07$, and no
symmetric input reaches the singlet. The default seed is now the \emph{uniform} metric
$l^2=1$, which is manifestly $g$-invariant; on it \eqref{eq:stokes} is exact
($\abs{\sigmacol_R}\sim 10^{-14}$ for a symmetric neutral input). The mutual-information
(van Raamsdonk) seed, \texttt{setEntangledMetric} --- entanglement-derived edge
lengths that are constant on $g$-orbits, hence $g$-invariant --- works equally well.

Crucially, the singlet is read from the \emph{transport} at the symmetric seed, not
from a relaxation: the intertwining of \S\ref{sec:theorem} holds already at iteration
zero. Relaxation still runs and is well-behaved ($60$ iterations reduce the
stationarity residual from $204$ to $8.4$), but it is not what produces the singlet.

\subsection{(ii) The \texorpdfstring{$\omega$}{omega}-representation input}
The ``natural symmetric quark input'' is \emph{not} $[1,\omega,\omega^2]$ placed
identically on $A,B,C$ (that gives only $\sim 0.48$). It is the input that transforms
in the $\omega$-representation of $g$:
\begin{equation}
  \Pin\,\phi = \omega\,\phi, \qquad \omega = e^{2\pi i/3}.
  \label{eq:omegainput}
\end{equation}
Physically: the three quark colors carry the correct relative $\Z_3$ phase under the
geometric cycling, not the same vector repeated. The input representation is fixed by
$g$ --- a property of the construction --- \emph{before} any output is computed; the
theorem's content (\S\ref{sec:theorem}) is that the transport carries it to the
singlet line rather than scrambling it, as the failed rows of \S\ref{sec:results}
show it does for non-$\omega$-rep inputs.

\subsection{(iii) A \texorpdfstring{$g$}{g}-coherent signed read-out}
The period of a $1$-cochain around a hole triangle is orientation-dependent. The
read-out must use orientations coherent under $g$: the period of each hole is taken in
the orientation propagated from a chosen base hole by $g$, i.e.\ multiplied by the
orientation parity $\varepsilon$ of $g$ on it (the corresponding entry of $P$).
Equivalently, on the three holes of $R$ the read-out signs are the entries of $\Pout$.
The default per-window covector is \emph{not} $g$-coherent --- its signs are
inconsistent across windows ($[-1,-1,1],[1,-1,-1],[1,-1,1]$) --- and with raw sorted
periods the symmetry is broken in the read-out and the singlet is missed.

\section{The intertwining theorem}
\label{sec:theorem}

Let $M:\C^9\to\C^3$ be the input$\to$result transport: $M$ sends the nine pinned
input-hole periods to the three result-hole periods of the carried harmonic. $M$ is
linear because the carried-representative map is linear.

\begin{lemma}[Equivariance on the base surface]
\label{lem:equivariant}
On the base surface (the holed geodesic sphere), $g$ is a metric-preserving simplicial
automorphism. It acts on hole-period vectors by a signed permutation $P$, the sign on a
hole being the orientation parity of $g$ on it; $P=\Pin\oplus\Pout$ because $g$ maps
input holes to input holes and result holes to result holes. The base-surface carried
map commutes with $g$:
\begin{equation}
  M\,\Pin = \Pout\,M .
  \label{eq:intertwine}
\end{equation}
\end{lemma}

\begin{proof}
$g$ preserves edge lengths, so it preserves the discrete Hodge inner products and
commutes with the coboundary and Hodge Laplacian; the least-norm harmonic realizing a
set of pinned periods is carried to the harmonic realizing the $g$-transported
periods. A hole $(a,b,c)$ (oriented, sorted) has period $\psi_{ab}+\psi_{bc}-\psi_{ac}$,
invariant under cyclic relabelling and negated under a transposition; $g$ sends
$(a,b,c)\mapsto(ga,gb,gc)$, whose sorted form differs by a permutation of parity
$\varepsilon\in\{\pm1\}$, so the period maps with sign $\varepsilon$. Collecting the
signs gives $P=\Pin\oplus\Pout$ and \eqref{eq:intertwine}.
\end{proof}

\begin{theorem}[The $\omega$-input transports to the singlet]
\label{thm:singlet}
Assume the exact equivariance \eqref{eq:intertwine}. Then $\Pout$ is a signed
$3$-cycle with eigenvalues $\{1,\omega,\omega^2\}$, and its $\omega$-eigenspace is the
one-dimensional color singlet line $\C\,[1,\omega,\omega^2]$ in $R$'s hole basis. For
any input $\phi$ with $\Pin\phi=\omega\phi$, the result $M\phi$ lies in that singlet
line. Moreover $M$ is rank $3$ with $\phi\notin\ker M$ (so $M\phi\neq 0$), and the
$\omega$-eigenspace of $\Pin$ is three-dimensional, so a three-parameter family of
symmetric inputs all transport to the singlet.
\end{theorem}

\begin{proof}
$g$ cycles $R$'s three holes, so $\Pout$ is a signed $3$-cycle; the product of its
three signs is $+1$ (verified: the eigenvalue angles are $\{0,\pm2\pi/3\}$), hence its
eigenvalues are the cube roots of unity, each eigenspace one-dimensional. The
$\omega$-eigenvector is the discrete Fourier mode $[1,\omega,\omega^2]$ --- the color
singlet. For $\Pin\phi=\omega\phi$, equation \eqref{eq:intertwine} gives
$\Pout(M\phi)=M(\Pin\phi)=\omega\,(M\phi)$, so $M\phi$ is an $\omega$-eigenvector of
$\Pout$, i.e.\ $M\phi\in\C\,[1,\omega,\omega^2]$; and $M\phi\neq0$ since $M$ is rank
$3$ and $\phi$ is not in its kernel. Finally $\Pin$ is three disjoint signed $3$-cycles
(each an $A$-hole, its $B$-image, its $C$-image), each with sign-product $+1$, so
$\omega$ occurs with multiplicity three.
\end{proof}

\begin{remark}[The measured object is equivariant only up to a triangulation residual]
\label{rem:residual}
$g$ is a metric-preserving automorphism of the base \emph{surface}, but the cobordism
is the base $\times I$, and the prism tetrahedralization's diagonal choices are not
$g$-invariant --- $g$ permutes vertices but not simplices. Hence the carried map
\emph{through} the $3$-complex is only approximately equivariant,
$\norm{M\Pin-\Pout M}/\norm{M}\approx 4.3\times10^{-2}$ (\S\ref{sec:results}), and the
$\omega$-input transports to the singlet line up to that residual --- overlap $0.999$
rather than exactly $1$. The residual is metric-independent (\S\ref{sec:results}); a
$g$-symmetric subdivision of the extrusion would remove it.
\end{remark}

\begin{remark}[Manifest $\Z_3$, and the relation to $S_3$]
\label{rem:z3}
What is geometrically realized is the color $\Z_3$: $g$ is its generator, the transport
intertwines it, and the singlet is its $\omega$-Fourier mode --- the $\Z_3$ mode that
survives restricting the totally antisymmetric $\varepsilon_{ijk}$ to the cyclic
subgroup. The full color Weyl group is $S_3$; its $\Z_3$ cyclic part is the one made
manifest by the tetrahedral geometry, and the conjugate symmetric input
($\Pin\phi=\omega^2\phi$) lands on the conjugate singlet $[1,\omega^2,\omega]$. The
singlet is thus not a tuned target but the forced image of an entire eigenspace --- the
sense in which it is emergent rather than imposed.
\end{remark}

\section{Numerical results}
\label{sec:results}

All figures are measured through the production C++ \texttt{TransportCobordism} with
the symmetric \texttt{TripartiteRegisterTopology}; the symmetry $g$ and the signed reps
$\Pin,\Pout$ are reconstructed from the icosahedral generators applied to the
cobordism's \emph{own} window labels (never Python-reconstructed holes). The figure of
merit is the normalized overlap
\begin{equation}
  \text{overlap} = \frac{\abs{\langle v_R,\,s\rangle}}{\norm{v_R}\,\norm{s}},
  \qquad s=[1,\omega,\omega^2],
  \label{eq:overlap}
\end{equation}
where $v_R$ is the carried result-hole period vector (read from $R$ after the carry,
never pinned).

\begin{center}
\setlength{\tabcolsep}{5pt}
\begin{tabular}{@{}p{6.6cm}cc@{}}
\toprule
Placement / metric / input & overlap with $s$ & note\\
\midrule
greedy windows, jitter, natural neutral input        & $\sim 0.74$ & windows inequivalent\\
symmetric windows, uniform, $[1,\omega,\omega^2]$ identical on $A,B,C$ & $0.48$ & wrong input (not $\omega$-rep)\\
symmetric windows, uniform, naive neutral, raw read  & $0.65$ & wrong (non-coherent) read-out\\
\textbf{symmetric windows, uniform, $\omega$-rep input} & $\mathbf{0.999}$ & the proton\\
\bottomrule
\end{tabular}
\end{center}

\noindent The structural and symmetry measurements:
\begin{itemize}[nosep]
  \item topology: $b_1=11$, dual complex valid (a manifold, no weld); metric uniform.
  \item $\Pout$ eigenvalues $\{1,\omega,\omega^2\}$ (angles $\{0,\pm2\pi/3\}$): the color $\Z_3$.
  \item intertwining $\norm{M\Pin-\Pout M}/\norm{M} = 4.26\times10^{-2}$.
  \item the three $\omega$-eigenvectors of $\Pin$ map to the singlet with overlaps
        $0.9991,\,0.9998,\,0.9991$; the small spread is the fixed-resolution
        triangulation artifact (basis-dependent), with $0.999$ the reported floor.
  \item charge conservation \eqref{eq:stokes} exact ($10^{-14}$) for a symmetric
        neutral input on the uniform metric; the singlet is traceless,
        $1+\omega+\omega^2=0$.
\end{itemize}

\paragraph{The residual is a triangulation artifact, not symmetry breaking.} The
intertwining is not exactly zero, and the cause is \emph{not} the metric: replacing the
uniform seed by a $g$-invariant \emph{non-degenerate} metric (edge lengths constant on
$g$-orbits) gives the same residual, $4.19\times10^{-2}$ versus $4.26\times10^{-2}$.
The base surface, its holes, and the metric are all $g$-symmetric; what is not is the
prism tetrahedralization (the $\times I$ extrusion's diagonal choices), so the discrete
carried map is only approximately equivariant at this fixed resolution
(Remark~\ref{rem:residual}). The $0.999$ overlap, not the exact $1$, is the signature
of that artifact; a $g$-symmetric subdivision would remove it. This is the conclusive
figure that the symmetry is respected. \emph{That subdivision has since been built}
(the symmetric apex interior, \#413): replacing the prism extrusion with a
label-independent stacking --- each base triangle coning to a face-apex, the gap
octahedra split along the canonical dual edge $f_1 f_2$, so no vertex-label sort chooses
a diagonal --- drives the residual from $4.26\times10^{-2}$ to $4.5\times10^{-14}$
(machine zero) and the singlet overlap to $1.000000$, confirming the artifact diagnosis
exactly.

\section{Implementation}
\label{sec:impl}

The change is confined to \texttt{TripartiteRegisterTopology}.
\begin{itemize}[nosep]
  \item \texttt{geodesicTwoSphere} now returns its edge-midpoint map alongside the
        faces, so the window generator can lift the $C_3$ permutations.
  \item \texttt{symmetricWindows} generates the four windows from the four $C_3$
        generators, a seed corner per window, and the midpoint lift, with the
        build-time assertions of Proposition~\ref{prop:windows}; it replaces the greedy
        twelve-hole pick.
  \item the default metric seed is uniform $l^2=1$ (the jitter seed is removed); the van
        Raamsdonk seed (\texttt{setEntangledMetric}) remains as a $g$-invariant
        alternative.
\end{itemize}
The result block $R$ still \emph{emerges} --- it is never pinned or hand-placed: the
read-out pins the three inputs on $A,B,C$ and returns $R$ to be read off the transport
after the carry. Validation reconstructs $g$ and $\Pin,\Pout$ from the cobordism's own
labels and checks Theorem~\ref{thm:singlet} end-to-end. A worked demonstration is
\nolinkurl{examples/cobordism/proton_tripartite_junction.py}.

\section{Discussion}
\label{sec:discussion}

The proton singlet is now \emph{forced by symmetry}. With symmetric windows and a
symmetric metric the transport intertwines the color $\Z_3$; the color-symmetric quark
input is, by definition, the $\omega$-representation; and any such input lands on the
unique $\omega$-eigenvector of the result, the singlet $[1,\omega,\omega^2]$. Nothing
about the singlet is inserted: it is the image of an entire eigenspace, and a
three-parameter family of symmetric inputs all reach it. The earlier $0.74$ ceiling was
not a limitation of the geometry but of \emph{inequivalent} windows; the earlier $0.48$
and $0.65$ ``failures'' were the wrong input and the wrong read-out, not the wrong
construction.

Two honest caveats. First, the geometrically realized symmetry is $\Z_3$, the cyclic
subgroup of the color Weyl group $S_3$; the cyclic singlet is what is forced. Second,
the intertwining is exact only up to a fixed-resolution triangulation artifact
($\sim4\times10^{-2}$), so the measured overlap is $0.999$ rather than $1$
(Remark~\ref{rem:residual}).

The path forward (\#380) reads emergent observables off the relaxed singlet geometry
--- net charge, a radius, the mass ratio --- and follows the photon (a worldline whose
$l^2$ relaxes through null) into re-absorption.

\section{Conclusion}
The trivalent junction with \emph{symmetric} register windows resolves the last gap in
building the proton from quarks. Placing the four windows as one tetrahedral $A_4$ orbit
makes them equivalent; on a symmetry-respecting metric the input$\to$result transport
intertwines the color $\Z_3$; and the natural color-symmetric input then transports to
the color singlet with overlap $0.999$. Symmetric windows are necessary but not
sufficient --- the metric, the input representation, and the read-out must all respect
the same symmetry --- and when they do, the proton's color singlet is not assembled but
\emph{compelled}.

\appendix
\section{Explicit data}
\label{app:data}

The four $C_3$ generators as $12$-vertex permutations (image of vertex $i$ at position
$i$), each cycling its window's vertex block:
\[
\begin{array}{ll}
\winA\ (2\!\to\!8\!\to\!10): & [4,3,8,9,5,0,7,11,10,1,2,6]\\
\winB\ (1\!\to\!4\!\to\!7):  & [3,4,0,2,7,8,5,1,6,11,9,10]\\
\winC\ (0\!\to\!6\!\to\!9):  & [6,10,11,7,2,1,9,8,3,0,5,4]\\
\winR\ (3\!\to\!5\!\to\!11): & [10,6,1,5,9,11,2,0,4,8,7,3]
\end{array}
\]
The seed corner of window $w$ is $(\,v,\,m(v,n_1),\,m(v,n_2)\,)$ with $(v;n_1,n_2)$
equal to
\[
  (2;0,1),\allowbreak\quad (1;6,10),\allowbreak\quad (0;3,4),\allowbreak\quad (3;2,7)
\]
respectively; the window is the $C_3$ orbit of its seed. The window-cycling element $g$
of Definition~\ref{def:g} is the $A_4$ rotation whose induced window permutation is
$\winA\to\winB\to\winC\to\winA$ with $\winR$ fixed; it acts on $R$'s three holes as a
$3$-cycle, giving $\Pout$ its eigenvalues $\{1,\omega,\omega^2\}$.

\end{document}
